\begin{abstract}
Tool-agent recovery methods combine several changes at once: additional
failure-conditioned rollouts, different advantage groups, recovery-specific
rewards, and worst-group reweighting.  A gain under such a bundle does not
identify which recovery mechanism caused it.  We conduct a controlled audit
of Minimax Recovery Policy Optimization (\method{}) for a Qwen3.5-9B LoRA
agent in executable, stateful tool environments.  Before training, we
calibrate task difficulty, freeze task--fault manifests and state-based
verifiers, and preregister three seed-level directional hypotheses.  We then
compare \method{} with GRPO, no-fission, and no-minimax arms while matching
primary rollout count and advantage grouping, fixing the branch budget for
fission-bearing arms, and pairing evaluation episodes.  Across seeds 23, 29,
and 43, none of the three hypothesized
mechanism advantages replicates: the pooled paired success differences are
0/96 for \method{} versus GRPO (95\% bootstrap CI
$[-0.0938,0.0938]$), $-2$/96 versus no-fission
($[-0.1042,0.0625]$), and $-14$/96 versus no-minimax
($[-0.2604,-0.0312]$).  The preregistered regression gate passes, so the
result is not explained by a catastrophic seed, but minimax weighting is
consistently harmful in the primary suite.  These findings do not reproduce
or contradict the full Fission-GRPO system; they show that, in our controlled
regime, branching and minimax mechanisms do not retain an identifiable
advantage after estimator matching.  The main contribution is therefore an
auditable protocol and a multi-seed negative result that narrows when
structured recovery objectives should be expected to help.
\end{abstract}
